\documentclass[9pt,twocolumn]{pnas-new}
\articletype{RESEARCH ARTICLE}

\templatetype{pnasresearcharticle}
% Keep normal two-column flow, but place Significance at the bottom-right on page 1.
\makeatletter
\additionalelement{%
\afterpage{%
\begin{sigstatements}
\sffamily
\mdfdefinestyle{pnassigstylebottom}{backgroundcolor=pnasblueback,linecolor=pnasblueback,fontcolor=pnasbluetext,innertopmargin=9.3pt,innerrightmargin=11.4pt,innerbottommargin=12pt,innerleftmargin=11.4pt,linewidth=0pt,userdefinedwidth=\columnwidth}
\@ifundefined{@significancestatement}{}{%
\begin{mdframed}[style=pnassigstylebottom]
\section*{\textcolor{pantone}{Significance}}
\raggedright\fontsize{8.5}{12}\selectfont\textcolor{black}\@significancestatement
\end{mdframed}}
\end{sigstatements}
}}
\makeatother
\renewcommand{\abscontent}{
\noindent{\absfont \theabstract\par}
\vskip10pt
\noindent{\keywordsfont \@ifundefined{@keywords}{}{\@keywords}\par}
\vskip12pt
}

\begin{document}

\title{Modeling Disease Networks using SIR Simulation}

\author[a]{E. Cruz-Martinez}
\author[a]{N. Delingat}
\author[a]{T. Eapen}
\author[a,1]{L. Fiechter}
\author[a]{J. Levi}

\affil[a]{Department of Mathematics, University of California, Los Angeles, CA 90095}

\leadauthor{Cruz-Martinez}

\significancestatement{This study provides insights into how the structure of social networks influences the spread of infectious diseases and highlights the critical need for targeted immunization strategies to efficiently suppress outbreaks. By utilizing a network-based framework and spectral analysis, the research identifies the disproportionate impact of highly connected and vulnerable individuals in transmission dynamics.}

\authorcontributions{Author contributions: }
\authordeclaration{The authors declare no competing interest}
\correspondingauthor{\textsuperscript{1}To whom correspondence should be addressed.\\ E-mail: lucfiechter@g.ucla.edu}

\keywords{Disease Networks $|$ SIR Simulation $|$ Epidemic Modeling $|$ Targeted Immunization}

\begin{abstract}
Understanding how infectious diseases spread through human populations remains a central challenge in public health, particularly in highly connected social networks. This project investigates disease propagation using a network-based modeling framework, with the primary research question: which group of individuals is most susceptible to infection, what is the minimum fraction of this group that must be immunized to significantly reduce or suppress disease spread, and what strategies can we use to suppress infection and sickness spread the fastest? Identifying targeted vaccination strategies is especially important because widespread immunization can be costly, slow, or logistically difficult, while poorly targeted interventions may fail to prevent outbreaks. The study models a social network derived from online or real-world interaction data, where individuals are represented as nodes and social interactions as weighted edges. Edge weights reflect contact frequency or transmission likelihood. Additional attributes such as disease contagiousness, recovery or mortality rates, and individual vulnerability are incorporated into the model. Using this structure, we simulate disease transmission using random-walk--based spreading and diffusion over the graph. Spectral properties of the network, particularly those derived from the graph Laplacian, are analyzed to understand connectivity, cluster structure, and the potential speed and reach of outbreaks. Through repeated simulation and analysis, we expect to identify that highly connected individuals (high-degree nodes), along with those possessing higher intrinsic vulnerability to infection, play a disproportionately important role in causing virus transmission. Targeted immunization of a relatively small subset of these individuals may dramatically reduce overall spread.
\end{abstract}

\dates{This manuscript was compiled on \today}
\doi{\url{www.pnas.org/cgi/doi/10.1073/pnas.XXXXXXXXXX}}

\maketitle
\thispagestyle{firststyle}
\ifthenelse{\boolean{shortarticle}}{\ifthenelse{\boolean{singlecolumn}}{\abscontentformatted}{\abscontent}}{}
\section*{Introduction}
We study disease spread on a contact network under a fixed vaccination budget. Individuals are represented as nodes, interactions as edges, and transmission is modeled with a stochastic SIR process. The central question is: for a given $k$, which immunized set $U$ best reduces outbreak severity? We evaluate practical immunization strategies by comparing two outcomes from simulation: peak infections $I_{\max}(U)$ and total infected $T_{\mathrm{inf}}(U)$. This framing directly matches the optimization problem in the following section and allows a consistent comparison of strategy performance as vaccination coverage changes.

\section*{Expected Results}
We expect the simulations to demonstrate that disease spread is not homogeneous across a real world graph, with infection risk among specific subgroups of individuals rather than being uniformly distributed. In particular, nodes with high degree and high weighted connectivity are anticipated to exhibit higher infection rates. Spectral analysis of the graph Laplacian is expected to reveal that tightly connected clusters facilitate rapid local spread, while bridge nodes connecting clusters control for global diffusion. Specifically networks with smaller spectral gaps are expected to show faster outbreaks and broader reach, whereas networks with pronounced community structure may exhibit localized outbreaks that take longer to propagate. Simulation results are anticipated to indicate that random immunization strategies require a relatively large fraction of the population to be vaccinated to achieve meaningful suppression of disease spread. In contrast, targeted immunization that lowers the spectral gap is expected to produce a disproportionate reduction in outbreak size and transmission speed.

\section*{Mathematical Problem Statement}

Let $G=(V,E)$ be a contact network with $|V|=n$ individuals. We run a discrete-time stochastic SIR process on $G$ with transmission rate $\beta$ and recovery rate $\gamma$.

Each node $v\in V$ has state $X_v(t)\in\{\mathrm{S},\mathrm{I},\mathrm{R}\}$ at time $t$. For each time step:
\begin{itemize}
    \item If $X_v(t)=\mathrm{S}$, then $v$ becomes infected with probability
    \[
    p_v(t)=1-\exp\!\left(-\beta\sum_{u\in N(v)} w_{uv}\,\mathbf{1}\{X_u(t)=\mathrm{I}\}\right),
    \]
    where $w_{uv}$ is the contact weight on edge $(u,v)$.
    \item If $X_v(t)=\mathrm{I}$, then $v$ recovers with probability $\gamma$.
    \item Recovered nodes remain in state $\mathrm{R}$, with no chance of reinfection.
\end{itemize}

Before the epidemic starts, choose an immunized set $U\subseteq V$ with budget
\[
|U|=k, \qquad 0\le k\le n,
\]
and initialize all $v\in U$ in state $\mathrm{R}$.

For a fixed immunization set $U$, let $I_t(U)$ be the number of infected nodes at time $t$, and define
\[
I_{\max}(U)=\max_t I_t(U),
\]
\[
T_{\mathrm{inf}}(U)=\left|\{v\in V:\ v\text{ becomes infected at least once}\}\right|.
\]

The optimization goal is to choose $U$ to minimize either expected peak infections
\[
\min_{U\subseteq V,\ |U|=k} \ \mathbb{E}[I_{\max}(U)],
\]
or expected outbreak size
\[
\min_{U\subseteq V,\ |U|=k} \ \mathbb{E}[T_{\mathrm{inf}}(U)].
\]

In practice the evaluation of these problems is not computationally tractable; instead we empirically analyze various simpler strategies in the construction of $U$.

\section*{Limitations}
Our current experiments are bounded by computational power. In particular we are only able to run our experiments on networks of a few thousand nodes, however the results may vary for out-of-distribution networks. Notably betweenness centrality is also not included because we are unable to perform this analysis our network data. 

\section*{Materials and Methods}
This project uses publicly available network data sources and established mathematical modeling frameworks to ensure reproducibility and methodological rigor. This will be broken down in the following subsections: 

\begin{itemize}
    \item \textbf{High-Resolution Synthetic Contact networks} 
    \item \textbf{Empirical Social Interaction Structure} 
    \item \textbf{Established Propagation Models} 
    \item \textbf{Centrality-Based Intervention Strategies} 
\end{itemize}

\subsection*{1. High-Resolution Synthetic Contact Networks}
We utilize synthetic human contact network datasets representing individual-level interactions within households, schools, workplaces, and group quarters. These networks simulate realistic, multi-layered social contact structures and capture heterogeneous mixing patterns across hundreds to thousands of individuals per network. Such datasets are commonly used in epidemiological modeling to approximate real-world transmission environments.

\subsection*{2. Empirical Social Interaction Structure}
The constructed person-to-person contact network is derived from shared environments (e.g., occupation, education, and household membership), allowing us to model transmission along socially meaningful interaction pathways rather than assuming homogeneous mixing.

\subsection*{3. Established Propagation Models}
Disease spread is modeled using a discrete-time stochastic SIR (Susceptible-Infected-Recovered) framework. The model captures the branching transmission dynamics, where infected nodes may infect multiple neighbors. It also accounts for competing exposure effects when multiple infected neighbors attempt to transmit infection to a single susceptible node.

\subsection*{4. Centrality-Based Intervention Strategies}
To evaluate immunization strategies, we employ graph-theoretic centrality measures to identify structurally influential nodes. These include:

\begin{itemize}
    \item \textbf{Degree centrality}: identifies locally high-contact ``hubs.''
    \item \textbf{Eigenvector centrality}: identifies globally influential nodes connected to other influential nodes.
\end{itemize}

These measures allow for the comparison of targeted vaccination strategies with random immunization. All computational methods rely on widely available open-source Python libraries such as \textbf{networkx} and \textbf{numpy}, ensuring the accessibility, reproducibility, and transparency of the modeling framework.

\dataavail{All data and code will be made available upon publication in an open-source repository.}

\acknow{We would like to acknowledge the support and resources provided during the course of this project.}

\showacknow{} % Display the acknowledgments section

% \bibsplit[3]
% Bibliography
% \bibliography{pnas-sample}

\end{document}
